\documentclass[12pt, a4paper]{article}

% ── Pacotes ───────────────────────────────────────────────────────────────────
\usepackage[utf8]{inputenc}
\usepackage[T1]{fontenc}
\usepackage[brazil]{babel}
\usepackage{geometry}
\usepackage{graphicx}
\usepackage{booktabs}
\usepackage{tabularx}
\usepackage{multirow}
\usepackage{array}
\usepackage{xcolor}
\usepackage{colortbl}
\usepackage{hyperref}
\usepackage{amsmath}
\usepackage{listings}
\usepackage{enumitem}
\usepackage{float}
\usepackage{caption}
\usepackage{subcaption}
\usepackage{setspace}
\usepackage{fancyhdr}
\usepackage{titlesec}
\usepackage{mdframed}
\usepackage{tcolorbox}
\tcbuselibrary{most}

\geometry{
  top=2.5cm, bottom=2.5cm,
  left=3cm, right=2cm
}

% ── Cores ─────────────────────────────────────────────────────────────────────
\definecolor{ufpiblue}{RGB}{0, 51, 102}
\definecolor{lightblue}{RGB}{230, 240, 255}
\definecolor{lightgreen}{RGB}{220, 240, 220}
\definecolor{lightred}{RGB}{255, 230, 230}
\definecolor{lightyellow}{RGB}{255, 250, 220}
\definecolor{codegray}{RGB}{245, 245, 245}
\definecolor{successgreen}{RGB}{0, 128, 0}
\definecolor{warningred}{RGB}{180, 0, 0}

% ── Estilo de código ──────────────────────────────────────────────────────────
\lstset{
  backgroundcolor=\color{codegray},
  basicstyle=\ttfamily\footnotesize,
  breaklines=true,
  frame=single,
  rulecolor=\color{gray!40},
  xleftmargin=1em,
  xrightmargin=1em,
  aboveskip=0.5em,
  belowskip=0.5em
}

% ── Cabeçalho/rodapé ──────────────────────────────────────────────────────────
\pagestyle{fancy}
\fancyhf{}
\rhead{\textcolor{gray}{\footnotesize Questão 4 — Destilação de Conhecimento}}
\lhead{\textcolor{gray}{\footnotesize UFPI -- Tópicos em IA -- 2026.1}}
\rfoot{\thepage}
\renewcommand{\headrulewidth}{0.4pt}

% ── Título de seções ──────────────────────────────────────────────────────────
\titleformat{\section}{\large\bfseries\color{ufpiblue}}{}{0em}{\thesection\quad}[\titlerule]
\titleformat{\subsection}{\normalsize\bfseries\color{ufpiblue!80}}{}{0em}{\thesubsection\quad}

% ── Caixas coloridas ──────────────────────────────────────────────────────────
\newtcolorbox{destaquebox}[1][]{
  colback=lightblue, colframe=ufpiblue,
  fonttitle=\bfseries\small, title=#1,
  left=6pt, right=6pt, top=4pt, bottom=4pt
}
\newtcolorbox{successbox}[1][]{
  colback=lightgreen, colframe=successgreen!70,
  fonttitle=\bfseries\small, title=#1,
  left=6pt, right=6pt, top=4pt, bottom=4pt
}
\newtcolorbox{exemplobox}[1][]{
  colback=lightyellow, colframe=gray!50,
  fonttitle=\bfseries\footnotesize\color{ufpiblue}, title=#1,
  left=6pt, right=6pt, top=4pt, bottom=4pt,
  fontupper=\footnotesize
}
\newtcolorbox{warningbox}[1][]{
  colback=lightred, colframe=warningred!70,
  fonttitle=\bfseries\small, title=#1,
  left=6pt, right=6pt, top=4pt, bottom=4pt
}

% ═════════════════════════════════════════════════════════════════════════════
\begin{document}
% ═════════════════════════════════════════════════════════════════════════════

% ── Capa ──────────────────────────────────────────────────────────────────────
\begin{titlepage}
  \centering
  \vspace*{1cm}
  {\large\bfseries UNIVERSIDADE FEDERAL DO PIAUÍ}\\[0.3cm]
  {\normalsize Centro de Ciências da Natureza -- Departamento de Computação}\\[0.3cm]
  {\normalsize Tópicos em Inteligência Artificial -- Prof. Raimundo Moura -- 2026.1}\\
  \vspace{2cm}
  \rule{\linewidth}{0.5pt}\\[0.4cm]
  {\LARGE\bfseries\color{ufpiblue} Destilação de Conhecimento em\\[0.3cm]
  Modelos de Linguagem}\\[0.3cm]
  {\large Questão 4 do Trabalho Final}\\
  \rule{\linewidth}{0.5pt}\\
  \vspace{1.5cm}
  \begin{tabular}{ll}
    \textbf{Modelos}    & Qwen2.5-3B-Instruct (Professor) $\rightarrow$ Qwen2.5-1.5B-Instruct (Aluno) \\[0.3cm]
    \textbf{Técnica}    & Chain-of-Thought Distillation + LoRA (SFTTrainer) \\[0.3cm]
    \textbf{Dataset}    & DOMPI-2025 (Diários Oficiais dos Municípios do Piauí) \\[0.3cm]
    \textbf{Benchmark}  & 100 perguntas sobre documentos municipais piauienses \\
  \end{tabular}
  \vfill
  {\normalsize Teresina, \today}
\end{titlepage}

% ── Sumário ───────────────────────────────────────────────────────────────────
\tableofcontents
\newpage

% ═════════════════════════════════════════════════════════════════════════════
\section{Introdução}
% ═════════════════════════════════════════════════════════════════════════════


A destilação de conhecimento (\textit{Knowledge Distillation}), proposta originalmente
por \textit{Hinton et al.} (2015), é uma técnica que permite transferir o conhecimento
de um modelo maior (professor) para um modelo menor (aluno), tornando o aluno capaz de
aproximar o comportamento do professor com menos parâmetros e menor custo computacional.

\subsection{LLMs Usados para Destilação}

Modelos usados para a questão atual:

\begin{itemize}
  \item \textbf{Qwen2.5-3B $\rightarrow$ Qwen2.5-1.5B} -- escolha deste trabalho
\end{itemize}

A escolha de modelos da \textbf{mesma família} é recomendada pois compartilham o mesmo
vocabulário e tokenizador, facilitando a destilação por logits e reduzindo problemas de
incompatibilidade de distribuições.

% ═════════════════════════════════════════════════════════════════════════════
\section{Configuração Experimental}
% ═════════════════════════════════════════════════════════════════════════════

\subsection{Modelos}

\begin{table}[H]
\centering
\caption{Modelos utilizados no experimento de destilação}
\begin{tabular}{llrrr}
\toprule
\textbf{Papel} & \textbf{Modelo} & \textbf{Parâmetros} & \textbf{VRAM (fp16)} & \textbf{Quant.} \\
\midrule
Professor & Qwen/Qwen2.5-3B-Instruct & 3B & $\sim$6 GB & fp16 \\
Aluno     & Qwen/Qwen2.5-1.5B-Instruct & 1.5B & $\sim$3 GB & fp16 \\
\bottomrule
\end{tabular}
\end{table}

Ambos os modelos pertencem à família \textbf{Qwen2.5} (Alibaba Cloud), compartilhando
a mesma arquitetura Transformer com RoPE, SwiGLU e GQA. A versão \texttt{-Instruct}
foi escolhida por ter passado por ajuste fino supervisionado (SFT) e alinhamento por
RLHF, sendo adequada para tarefas de perguntas e respostas sobre documentos.

\subsection{Dataset Sintético}

O dataset de destilação foi gerado sinteticamente a partir do corpus
\textbf{DOMPI-2025}.

\textbf{Processo de geração:}
\begin{enumerate}
  \item Seleção de 100 trechos representativos do corpus DOMPI-2025
  \item Para cada trecho, o \textbf{professor} recebe um prompt solicitando raciocínio
        passo a passo (\textit{Chain-of-Thought}) e resposta final em formato JSON
  \item O resultado é salvo em \texttt{distillation\_dataset.jsonl} com os campos
        \texttt{instruction}, \texttt{input}, \texttt{output}, \texttt{reasoning} e \texttt{answer}
\end{enumerate}

\begin{exemplobox}[Exemplo de entrada do dataset sintético]
\textbf{Instrução:} Qual é o CEP da sede da Prefeitura de Aroazes?\\[0.3em]
\textbf{Documento:} ESTADO DO Piauí Prefeitura Municipal DE AROAZES C.N.P.J.:
06.554.984/0001-39 Av. 27 de Fevereiro, 691, Centro CEP: 64310-000\\[0.3em]
\textbf{Reasoning (professor):} To find the CEP of the Aroazes City Hall, I will
identify the relevant field in the document: `CEP: 64310-000'. The CEP is a
Brazilian postal code format (XXXXX-XXX). Therefore, the answer is 64310-000.\\[0.3em]
\textbf{Answer (professor):} 64310-000
\end{exemplobox}

\subsection{Técnica de Destilação: Chain-of-Thought Distillation}

A abordagem utilizada segue o paradigma de \textbf{destilação por Chain-of-Thought
(CoT)}, onde o aluno aprende não apenas as respostas finais do professor, mas também
o raciocínio intermediário. Isso é distinto da destilação clássica por KL-Divergence
de logits (Hinton et al., 2015) e se aproxima da abordagem de \textit{Instruction
Tuning} com dados gerados por LLMs maiores.

\subsection{Hiperparâmetros de Treinamento}

\begin{table}[H]
\centering
\caption{Hiperparâmetros utilizados no fine-tuning do aluno}
\begin{tabular}{ll}
\toprule
\textbf{Parâmetro} & \textbf{Valor} \\
\midrule
Método de fine-tuning   & LoRA (Low-Rank Adaptation) \\
Rank LoRA ($r$)         & 16 \\
Alpha LoRA              & 32 \\
Módulos alvo            & \texttt{q\_proj}, \texttt{v\_proj} \\
Dropout LoRA            & 0.05 \\
Épocas                  & 3 \\
Batch efetivo           & 16 (2 $\times$ 8 grad. accum.) \\
Learning rate           & $2 \times 10^{-4}$ \\
Scheduler               & Linear com warmup \\
Comprimento máximo      & 512 tokens \\
Precisão                & fp16 \\
\bottomrule
\end{tabular}
\end{table}

\subsection{Hardware Utilizado}

\begin{itemize}
  \item \textbf{GPU:} NVIDIA GeForce RTX 4050 (6 GB VRAM)
  \item \textbf{CPU:} Intel Core i5 13ª geração
  \item \textbf{RAM:} 16 GB
  \item \textbf{SO:} Windows 11 + Python 3.12
\end{itemize}

% ═════════════════════════════════════════════════════════════════════════════
\section{Metodologia de Avaliação}
% ═════════════════════════════════════════════════════════════════════════════

\subsection{Avaliação Qualitativa: Benchmark de 100 Perguntas}

O benchmark consiste em 100 pares \texttt{(instrução, contexto, resposta\_esperada)}
extraídos do DOMPI-2025, cobrindo tipos de pergunta como:
extração de CNPJ, endereços, valores de contrato, números de decreto, modalidades de
licitação e obrigações legais.

A avaliação de acerto usou correspondência textual normalizada:
\begin{enumerate}
  \item Normalização: minúsculas, remoção de pontuação e espaços extras
  \item Critério 1 -- \textit{Substring}: gabarito contido na resposta ou vice-versa
  \item Critério 2 -- \textit{Word Overlap}: $\geq 70\%$ das palavras do gabarito
        presentes na resposta
\end{enumerate}

\subsection{Avaliação Complementar: LLM-as-Judge}

Como complemento à métrica automática de \textit{substring}/\textit{word overlap},
as 100 respostas do aluno (antes e depois da destilação) foram também submetidas a
uma \textbf{avaliação por um LLM juiz} (\textit{LLM-as-Judge}), seguindo o protocolo
proposto por \textit{Zheng et al.} (2023).

Para cada uma das 100 perguntas, o juiz recebeu:
\begin{enumerate}
  \item A pergunta original (\texttt{instruction});
  \item A resposta de referência / gabarito (\texttt{output\_esperado});
  \item A resposta do modelo avaliado (\texttt{resposta\_aluno\_antes} ou
        \texttt{resposta\_aluno\_depois});
\end{enumerate}
e respondeu de forma binária -- \textbf{CORRETO} ou \textbf{INCORRETO} -- acompanhada
de uma justificativa curta, considerando a resposta correta mesmo quando havia
diferenças de formatação, ordem das palavras, pequenos erros ortográficos ou
informações adicionais irrelevantes, desde que o fato central solicitado estivesse
certo e sem contradições ou dados inventados (\textit{alucinações}).

Essa abordagem é mais robusta que a comparação por \textit{substring}, pois consegue
capturar:
\begin{itemize}
  \item Respostas \textbf{semanticamente corretas} mas com redação diferente do gabarito
        (que o \textit{substring matching} classificaria como erro);
  \item Respostas com o \textbf{trecho certo presente}, mas acompanhadas de
        \textbf{contradições internas} ou dados fabricados (alucinações), que o
        \textit{substring matching} classificaria erroneamente como acerto.
\end{itemize}

% ═════════════════════════════════════════════════════════════════════════════
\section{Resultados}
% ═════════════════════════════════════════════════════════════════════════════

\subsection{Métricas Quantitativas}

\begin{table}[H]
\centering
\caption{Comparação das métricas quantitativas antes e depois da destilação}
\begin{tabular}{lccc}
\toprule
\textbf{Modelo} & \textbf{Perplexidade} & \textbf{Entropia Cruzada} & \textbf{Acurácia Tokens} \\
\midrule
\rowcolor{lightblue}
Professor (3B)    & 4,50 & 1,5050 & 68,08\% \\
Aluno ANTES (1.5B) & 5,49 & 1,7022 & 65,65\% \\
\rowcolor{lightgreen}
Aluno DEPOIS (1.5B) & \textbf{4,23} & \textbf{1,4434} & \textbf{68,65\%} \\
\midrule
\multicolumn{4}{l}{\textit{Variação do aluno (antes $\rightarrow$ depois):}} \\
$\Delta$ & $-1,26\ (-22{,}9\%)$ & $-0,2588\ (-15{,}2\%)$ & $+2,98$ pp \\
\multicolumn{4}{l}{\textit{Gap aluno-depois vs. professor:}} \\
$\Delta$ & $-0,27$ & $-0,0616$ & $+0,57$ pp \\
\bottomrule
\end{tabular}
\end{table}

\begin{successbox}[Resultado-chave das métricas quantitativas]
O aluno destilado \textbf{superou o professor} em todas as três métricas:
PPL $4{,}23 < 4{,}50$, EC $1{,}4434 < 1{,}5050$ e Acurácia $68{,}65\% > 68{,}08\%$.
Isso indica que o fine-tuning especializou o aluno no domínio do dataset (DOMPI-2025),
tornando-o mais eficiente que o próprio professor nesse subconjunto.
\end{successbox}

\subsection{Benchmark Qualitativo: Acertos nas 100 Perguntas (Substring/Word Overlap)}

\begin{table}[H]
\centering
\caption{Resultado do benchmark qualitativo (100 perguntas) -- critério automático}
\begin{tabular}{lcc}
\toprule
\textbf{Métrica} & \textbf{Quantidade} & \textbf{\%} \\
\midrule
Professor acertou              & 55 & 55\% \\
Aluno ANTES acertou            & 36 & 36\% \\
\rowcolor{lightgreen}
Aluno DEPOIS acertou           & \textbf{42} & \textbf{42\%} \\
\midrule
Convergência (ambos acertaram) & 28 & 28\% \\
Aluno acertou, professor errou & 14 & 14\% \\
Professor acertou, aluno errou & 27 & 27\% \\
\midrule
\rowcolor{lightgreen}
Aluno melhorou após destilação & 16 & 16\% \\
\rowcolor{lightred}
Aluno piorou após destilação   & 10 & 10\% \\
\bottomrule
\end{tabular}
\end{table}

O aluno passou de \textbf{36\%} para \textbf{42\%} de acertos (+6 pontos percentuais),
melhorando em 16 questões com apenas 10 regressões. O gap em relação ao professor
(55\% vs 42\%) reflete a diferença de capacidade entre os modelos (3B vs 1.5B).

\subsection{Benchmark Qualitativo: Acertos nas 100 Perguntas (LLM-as-Judge)}

Para verificar a robustez do resultado acima, as mesmas 100 respostas do aluno
(antes e depois) foram reavaliadas por um LLM juiz, conforme descrito na
Seção~3.2. A Tabela~\ref{tab:llm-judge} resume o resultado.

\begin{table}[H]
\centering
\caption{Resultado do benchmark qualitativo (100 perguntas) -- avaliação por LLM-as-Judge}
\label{tab:llm-judge}
\begin{tabular}{lcc}
\toprule
\textbf{Métrica} & \textbf{Quantidade} & \textbf{\%} \\
\midrule
\rowcolor{lightgreen}
Aluno ANTES acertou (juiz)      & \textbf{69} & \textbf{69\%} \\
Aluno DEPOIS acertou (juiz)     & 62 & 62\% \\
\midrule
Convergência -- ambos corretos    & 52 & 52\% \\
Convergência -- ambos incorretos  & 21 & 21\% \\
Aluno melhorou após destilação   & 10 & 10\% \\
\rowcolor{lightred}
Aluno piorou após destilação     & \textbf{17} & \textbf{17\%} \\
\bottomrule
\end{tabular}
\end{table}

\begin{warningbox}[Resultado-chave da avaliação por LLM-as-Judge]
Ao contrário do que sugere a métrica automática de \textit{substring}, a avaliação
semântica por LLM-as-Judge indica que o aluno \textbf{piorou} após a destilação:
de \textbf{69\%} de acertos antes para \textbf{62\%} depois ($-7$ pontos percentuais),
com \textbf{17} questões que regrediram contra apenas \textbf{10} que melhoraram.
Em \textbf{52} das 100 perguntas o julgamento não mudou (ambas as respostas corretas)
e em \textbf{21} ambas seguiram incorretas.
\end{warningbox}

As principais causas identificadas pelo juiz para as regressões (piora) incluem:
\begin{itemize}
  \item \textbf{Alucinações introduzidas após o fine-tuning}: em vários casos a
        resposta ``depois'' inventa detalhes que não estão no documento-fonte
        (por exemplo, um endereço ``em frente ao Banco do Brasil'' inexistente no
        texto original, ou números de CNPJ/CEP diferentes do gabarito);
  \item \textbf{Perda de informação corretamente extraída antes}: em algumas
        perguntas o modelo \textit{antes} da destilação continha a resposta certa
        (ex.: número de CNPJ, matrícula, data), e o modelo \textit{depois} passou a
        negar ter a informação (``não há informações específicas\ldots'');
  \item \textbf{Contradições internas}: respostas que mencionam o valor correto mas
        o descrevem de forma incoerente com o próprio texto (ex.: informar
        ``R\$~25.000,00'' e, na mesma frase, escrever por extenso ``vinte e quatro
        mil reais'').
\end{itemize}

Por outro lado, o juiz também identificou \textbf{10 melhoras genuínas} após a
destilação, em que a resposta ``depois'' extraiu corretamente uma informação que a
resposta ``antes'' não continha ou apresentava de forma incompleta/confusa (por
exemplo, a matrícula do servidor foi corrigida de ``1234-'' para ``1234-5'', e o
percentual de desconto do IPTU, ausente na resposta ``antes'', apareceu correto e
completo na resposta ``depois'').

\subsection{Comparação entre os Dois Métodos de Avaliação}

\begin{table}[H]
\centering
\caption{Comparação: métrica automática (substring/overlap) vs. LLM-as-Judge}
\begin{tabular}{lcc}
\toprule
\textbf{Método de avaliação} & \textbf{Aluno ANTES} & \textbf{Aluno DEPOIS} \\
\midrule
Substring / Word Overlap & 36\% & 42\% ($+6$ pp) \\
LLM-as-Judge             & 69\% & 62\% ($-7$ pp) \\
\bottomrule
\end{tabular}
\end{table}

Os dois métodos divergem tanto nos valores absolutos quanto na \textbf{direção} do
efeito da destilação. A métrica de \textit{substring} é mais rígida quanto à forma
exata da resposta e por isso subestima o desempenho geral (ambos os modelos ficam
abaixo de 45\%), mas identifica uma melhora líquida após o fine-tuning. Já o
LLM-as-Judge, por avaliar o \textbf{conteúdo semântico} da resposta, reconhece um
percentual bem maior de acertos em ambos os momentos, porém aponta uma
\textbf{piora líquida} após a destilação -- sugerindo que parte do ganho medido pela
métrica automática pode estar relacionado a respostas ``depois'' mais curtas e
diretas (que casam melhor com o gabarito em nível de \textit{string}), mas que em
alguns casos vêm acompanhadas de alucinações que um avaliador semântico consegue
detectar e a métrica automática não.

% ═════════════════════════════════════════════════════════════════════════════
\section{Discussão}
% ═════════════════════════════════════════════════════════════════════════════

\subsection{Houve Transferência de Conhecimento?}

\begin{destaquebox}[Resposta objetiva ao enunciado]
\textbf{Sim, houve transferência de conhecimento}, mas com evidências mistas sobre a
sua qualidade. As três métricas quantitativas confirmam a transferência de forma
inequívoca:
\begin{itemize}
  \item Perplexidade reduziu \textbf{22,9\%} (5,49 $\rightarrow$ 4,23)
  \item Entropia cruzada reduziu \textbf{15,2\%} (1,7022 $\rightarrow$ 1,4434)
  \item Acurácia de tokens aumentou \textbf{+2,98 pp} (65,65\% $\rightarrow$ 68,65\%)
  \item O aluno destilado \textbf{superou o professor} no domínio avaliado
\end{itemize}
Entretanto, ao nível de \textbf{corretude semântica das respostas} (LLM-as-Judge), o
aluno destilado \textbf{piorou} (69\% $\rightarrow$ 62\%), enquanto pela métrica de
\textit{substring} ele \textbf{melhorou} (36\% $\rightarrow$ 42\%). Isso mostra que
menor perplexidade/entropia cruzada não implica necessariamente respostas mais
corretas ou mais confiáveis do ponto de vista factual.
\end{destaquebox}

\subsection{Por que o Aluno Superou o Professor nas Métricas?}

Um resultado aparentemente contra-intuitivo é o aluno destilado ter superado o
professor nas métricas quantitativas (PPL $4{,}23 < 4{,}50$). Isso ocorre porque:

\begin{enumerate}
  \item As métricas são calculadas sobre o \textbf{mesmo dataset} usado no treino
        do aluno (\texttt{distillation\_dataset.jsonl}), gerando um viés de ajuste
        (\textit{in-distribution advantage})
  \item O fine-tuning com LoRA \textbf{especializou} o aluno no estilo de resposta
        do professor nesse domínio específico (documentos municipais piauienses)
  \item O professor, sendo um modelo genérico, não está especializado nesse formato
        de avaliação
\end{enumerate}

Em um benchmark \textit{out-of-distribution}, espera-se que o professor ainda
supere o aluno. Nos acertos do benchmark qualitativo, isso se confirma: professor 55\%
vs aluno 42\% (métrica \textit{substring}).

\subsection{Por que a Métrica Automática e o LLM-as-Judge Divergem?}

A divergência entre os dois métodos de avaliação qualitativa (Seção~4.4) é, em si,
um resultado relevante para este trabalho:

\begin{enumerate}
  \item A métrica de \textbf{substring/word overlap} recompensa respostas curtas e
        objetivas que contenham literalmente o gabarito, mas \textbf{não penaliza}
        alucinações ou contradições que apareçam ao lado da informação correta;
  \item O \textbf{LLM-as-Judge} avalia o \textbf{significado} da resposta como um
        todo e penaliza tanto a ausência da informação correta quanto a presença de
        \textbf{informações fabricadas} (mesmo quando o dado correto também aparece
        no texto);
  \item O fine-tuning por LoRA parece ter ensinado o aluno a produzir respostas em um
        \textbf{formato} mais parecido ao do professor/gabarito (o que ajuda na
        métrica de \textit{substring}), mas em alguns casos também introduziu
        \textbf{maior propensão a alucinar} detalhes não presentes no documento
        original, prejudicando a avaliação semântica.
\end{enumerate}

Isso reforça a recomendação, já presente na Seção~4.5, de tratar a métrica
\textit{substring} com cautela e priorizar avaliações semânticas (LLM-as-Judge ou
revisão humana) para conclusões definitivas sobre corretude.

\subsection{Limitações e Observações}

\begin{itemize}
  \item \textbf{Respostas em inglês do professor:} o Qwen2.5-3B-Instruct gerou
        diversas respostas em inglês apesar do prompt em português. Isso reduziu
        artificialmente os acertos do professor no benchmark (critério de substring
        não casa com tradução). Usar um prompt mais diretivo em português mitigaria
        esse problema.

  \item \textbf{Tamanho do dataset:} 100 exemplos é um dataset pequeno para
        destilação. Com 1.000+ exemplos (como na Questão 2), os resultados seriam
        mais robustos e generalizáveis.

  \item \textbf{Avaliação por substring:} a métrica de acerto é aproximada e pode
        ter falsos positivos/negativos, como evidenciado pela divergência com o
        LLM-as-Judge (Seção~4.4 e 5.3). Métricas como ROUGE-L ou avaliação por
        múltiplos LLMs juízes (\textit{ensemble of judges}) tornariam a conclusão
        ainda mais robusta.

  \item \textbf{Viés do próprio LLM-as-Judge:} o avaliador semântico também é um
        modelo de linguagem e pode ter seus próprios vieses (por exemplo, maior
        tolerância a paráfrases ou sensibilidade excessiva a pequenos detalhes
        divergentes). O ideal seria complementar esta análise com uma amostra
        revisada por humanos.
\end{itemize}

% ═════════════════════════════════════════════════════════════════════════════
\section{Conclusão}
% ═════════════════════════════════════════════════════════════════════════════

Este trabalho demonstrou a viabilidade da destilação de conhecimento por
\textit{Chain-of-Thought} entre modelos da família Qwen2.5, aplicada ao domínio
de documentos oficiais dos municípios do Piauí.

Os principais resultados foram:

\begin{itemize}
  \item \textbf{Transferência confirmada nas métricas quantitativas:} redução de
        22,9\% na perplexidade e aumento de 2,98 pp na acurácia de tokens após a
        destilação
  \item \textbf{Especialização no domínio:} o aluno destilado superou o professor
        nas métricas quantitativas no domínio avaliado
  \item \textbf{Resultado qualitativo misto e dependente do método de avaliação:}
        pela métrica automática de \textit{substring}, o aluno melhorou (36\%
        $\rightarrow$ 42\%, +6 pp); já pela avaliação semântica com LLM-as-Judge, o
        aluno piorou (69\% $\rightarrow$ 62\%, $-7$ pp), com 17 regressões contra
        10 melhoras -- em grande parte associadas a alucinações introduzidas após
        o fine-tuning
  \item \textbf{Eficiência:} o aluno (1,5B parâmetros) alcançou desempenho comparável
        ao professor (3B parâmetros) com metade dos parâmetros, via fine-tuning LoRA
        que treinou apenas $\approx$0,5\% dos parâmetros do modelo

\end{itemize}

A abordagem demonstra que modelos menores podem ser especializados em domínios
específicos por meio de destilação com dados sintéticos gerados por modelos maiores,
reduzindo custos de inferência sem sacrificar significativamente a qualidade das
respostas no domínio alvo. Ao mesmo tempo, a divergência entre a métrica automática
e a avaliação por LLM-as-Judge alerta para o risco de conclusões precipitadas sobre
"melhora" de um modelo destilado quando a avaliação se baseia apenas em correspondência
textual: ganhos de formato e concisão podem mascarar uma perda de confiabilidade
factual (alucinações), que só uma avaliação semântica (ou humana) é capaz de capturar.

% ═════════════════════════════════════════════════════════════════════════════
\section*{Referências}
% ═════════════════════════════════════════════════════════════════════════════

\begin{enumerate}[label={[\arabic*]}]
  \item \label{dompi2025} Portelaa, G. et al. (2025). \textit{DOMPI-2025: Dataset dos
        Diários Oficiais dos Municípios do Piauí}.
        \url{https://huggingface.co/datasets/gutoportelaa/DOMPI-2025}
  \item \label{llmjudge} Zheng, L. et al. (2023). \textit{Judging LLM-as-a-Judge with
        MT-Bench and Chatbot Arena}. arXiv:2306.05685.
\end{enumerate}

\end{document}